\section{Conclusions}\label{sec:conclusions}
We have obtained local tube-volume bounds for regular Pfaffian hypersurfaces, including unbounded ones, and used them to bound a geometric condition number for classifiers under uniform and Gaussian sampling. For one hidden sigmoid layer with rational first-layer weights and fixed lattice constant, the tail is bounded by $O_{n,L,m}(w^n/t)$. The underlying Gauss-map degree bound $\md(V)\le n!L^nw^n$ has optimal width exponent, although sharpness of the probability estimate remains open.

\subsection{Improved bounds for neural networks}\label{subsec:improved-nn}
For general sigmoid networks, the Pfaffian bound still contains $2^{h(h-1)/2}$, with $h$ the number of hidden units. Conjecture~\ref{conj:multi-layer-degree} asks for polynomial bounds in the widths at fixed input dimension, depth and lattice constant. A corresponding estimate for every section degree would also improve the multilayer tube and robustness bounds.

The sigmoid chain has structure beyond a general Pfaffian chain. A chain element in layer $i$ has derivative polynomials of degree at most $2i$, depending on that element and on earlier layers, but not on other elements of its own layer. Moreover, for an input $u\in\R^{n_{i-1}}$ and preactivation $z^i=A^iu+b^i$, the differential of the layer map is
\[
\diff{F^i}(u)=\diag(\sigma'(z^i_1),\dots,\sigma'(z^i_{n_i}))A^i,
\qquad \sigma'(z^i_j)>0.
\]
It is natural to ask whether these properties permit a more efficient Rolle-type zero count. The usual Khovanskii argument (see, for example, the detailed proof in~\cite{lotz2026khovanskii}) processes chain elements successively; controlling the degree growth when processing a whole layer remains a difficulty. We do not assert a uniform improvement for arbitrary monotone Pfaffian activations: any such statement must account for the complexity of the activation itself, including its chain length.

Even for one hidden layer, sharpening the section bounds between $\md_0(V)$ and $\md_{n-1}(V)$ could reduce the constants and the higher-order coefficients in the tube estimate. The lower-bound construction for the top Gauss degree does not settle these questions.

\subsection{Non-smooth Pfaffian sets}\label{subsec:nonsmooth}
The main theorem requires $0$ to be a regular value of the defining function. Although the top-score tie locus can be nonsmooth where regular pairwise equality hypersurfaces meet, the union bound still applies in that situation. For a nonzero shallow sigmoid output, Corollary~\ref{cor:single-layer-singular-linear} removes the regular-value assumption from the linear-radius tube bound by passage to nearby regular levels. This retains the $O_{n,L}(w^n)$ coefficient, but does not recover the tube-radius power corresponding to a higher real codimension.

In the algebraic setting, Basu and Lerario~\cite{basu2021hausdorff} extend tube bounds to singular varieties through controlled smooth approximations and Hausdorff limits. Their construction uses genericity properties of complex algebraic polar varieties; see their Remark 1.2 for the difficulty of extending this method to the Pfaffian setting.

Pfaffian chains are real analytic and admit local holomorphic extensions satisfying the same differential identities. Moreover, critical loci of projections of regular Pfaffian sets can be expressed using Pfaffian Jacobian minors, with effective degree bounds. These local facts do not supply the global complex-algebraic genericity and deformation properties used in the algebraic approximation argument. Obtaining suitable replacements with controlled Pfaffian formats is an open direction for tube bounds that retain the real codimension. Covering and slicing methods, such as those of~\cite{zhang2025covering}, offer another route when one is willing to use different dimension-dependent constants.
